\chapter{热力学势}\label{ch:b2:thermodynamic-potentials}

快速拉伸一根橡皮筋，再把它贴到唇边：它是温的；
让它收缩，它就是凉的。在它上面挂一个重物，并用
吹风机加热它：它\emph{缩短}了。没有哪种钢弹簧如此，而
其原因是：橡皮筋的张力并非能量问题，而是
熵——被拉伸的链排列的方式更少。
要分析这类问题，需要弄清恒温
或恒温恒压下的系统倾向于怎样变化，以及
能从中提取多少功。这类问题，内能和
熵只能笨拙地回答，因为它们说的是
孤立系统的语言。本章引入为实验室条件而构造的
两个函数，即\emph{亥姆霍兹自由能} $F = U -
TS$ 与\emph{吉布斯自由焓} $G = H - TS$，表明它们
度量可回收的功，并随趋于平衡而减小，并
由其微分导出麦克斯韦关系，它们联系着
原本无人会想到相关的量，并将整个理论用于
两相平衡，其中 $G$ 决定哪一相胜出，
并给出克拉珀龙公式。

\begin{omfigure}
\includegraphics[width=0.58\linewidth]{images/book4/ai/b4-28-ice-lake-dawn.jpg}

{\small 黎明时分湖面上冰与水共存：在熔
点时两相每千克的自由焓相同，而一
微小的压强或温度变化即可打破平衡——其
幅度由克拉珀龙斜率给出。}
\end{omfigure}

\section{基本恒等式及其变量}

\begin{proposition}[基本恒等式]\label{prop:b2:thermodynamic-potentials:identity}
对于组分固定的封闭系统，内能
视为熵与体积的函数，满足
\[
\dd U = T\,\dd S - P\,\dd V, \qquad
T = \Big(\frac{\partial U}{\partial S}\Big)_V, \quad
P = -\Big(\frac{\partial U}{\partial V}\Big)_S ,
\]
而焓 $H = U + PV$ 满足 $\dd H = T\,\dd S + V\,\dd P$。
每个函数在其\emph{自然变量}\index{自然变量}
——$(S,V)$ 对应 $U$，$(S,P)$ 对应 $H$——包含
系统的全部热力学：其偏导数给出其他
状态参量。热力学第二定律指出，孤立系统
（$U$、$V$ 固定）的平衡态是 $S$ 取最大值的状态。
\end{proposition}

\begin{proof}
可逆过程的热力学第一定律为 $\dd U = \delta Q_{\text{可逆}}
+ \delta W_{\text{可逆}} = T\,\dd S - P\,\dd V$，只涉及状态
函数，因此对平衡态之间的任一无穷小
变化均成立。加上 $\dd(PV)$ 即得 $H$ 的表达式。
\end{proof}

\begin{remark}[为何引入新函数]\label{rem:b2:thermodynamic-potentials:why}
实验并非在固定熵下进行。化学家的烧瓶处于
室温与大气压强之下；置于热源中的
活塞则工作于固定 $T$ 的条件下。
人们希望得到自然变量为 $(T,V)$ 与 $(T,P)$ 的函数，它们在这些
条件下所起的作用，恰如 $S$ 对孤立系统所起的：判断
系统向何方演化，以及它能输出多少功。一个
\emph{勒让德变换}——减去 $TS$，把 $S$ 换成 $T$ 作为
自变量，由此构造出它们。
\end{remark}

\section{自由能与自由焓}

\begin{definition}[亥姆霍兹自由能，吉布斯自由焓]\label{def:b2:thermodynamic-potentials:FG}
\[
F = U - TS \quad\text{(\emph{自由能}\index{自由能})}, \qquad
G = H - TS = U + PV - TS \quad\text{(\emph{自由焓}\index{自由焓})}.
\]
二者均为广延状态函数，以焦耳为单位。
\end{definition}

\begin{theorem}[微分与导出量]\label{thm:b2:thermodynamic-potentials:dFdG}
\[
\dd F = -S\,\dd T - P\,\dd V, \qquad
\dd G = -S\,\dd T + V\,\dd P ;
\]
于是，在 $(T,V)$ 为 $F$ 的自然变量、$(T,P)$ 为 $G$ 的自然变量时，
\[
S = -\Big(\frac{\partial F}{\partial T}\Big)_V, \quad
P = -\Big(\frac{\partial F}{\partial V}\Big)_T, \qquad
S = -\Big(\frac{\partial G}{\partial T}\Big)_P, \quad
V = \Big(\frac{\partial G}{\partial P}\Big)_T ,
\]
并由（吉布斯–亥姆霍兹）关系得 $U = F - T(\partial F/\partial T)_V =
-T^2\,\partial(F/T)/\partial T$，$H = -T^2\,\partial(G/T)/\partial T$。
只要知道 $F(T,V)$ 或 $G(T,P)$，便掌握了一切：状态
方程、熵、能量与热容。
\end{theorem}

\begin{proof}
$\dd F = \dd U - T\dd S - S\dd T = -S\dd T - P\dd V$；对 $G$ 同理，
由 $\dd H$ 出发。读出偏导数；对于吉布斯–亥姆霍兹关系，
$F - T\,\partial_TF = F + TS = U$，且 $\partial_T(F/T) = (T\partial_TF -
F)/T^2 = -U/T^2$。
\end{proof}

\begin{example}[理想气体]\label{ex:b2:thermodynamic-potentials:perfect}
对 $n$ 摩尔，$U = nc_{V}T + U_0$ 与 $S = nc_V\ln T + nR\ln V + S_0$
给出 $F(T,V) = -nRT\ln V + f(T)$：于是 $P = -\partial_VF = nRT/V$
——状态方程由 $F$ 导出——而 $S = -\partial_TF$
即得熵。在变量 $(T,P)$ 下，$G(T,P) = nRT\ln P +
g(T)$，于是 $V = \partial_PG = nRT/P$，并且\emph{摩尔}自由焓
$\mu(T,P) = \mu^\circ(T) + RT\ln(P/P^\circ)$——即\emph{化学
势}\index{化学势}，它随
压强增大：气体自高 $\mu$ 流向低，此处即由高压流向低
压强。
\end{example}

\section{最大功与演化方向}

\begin{theorem}[单热源过程]\label{thm:b2:thermodynamic-potentials:monothermal}
与温度为 $T_0$ 的单一热源交换热量的系统，其
初温与终温均为 $T_0$，在任一
过程中接受的功为
\[
W \ge \Delta F , \qquad\text{即}\qquad W_{\text{回收}} = -W \le -\Delta F :
\]
$F$ 的减少量，即定温下可从系统中回收的
\emph{最大可回收功}（在可逆过程中达到）——
因此称为“自由能”：$U$ 中可转化为功的部分，其
余下的 $TS$ 为束缚部分。若过程还是外压恒定的
（外部压强为 $P_0$，初、终态压强均为 $P_0$），功
\emph{除大气所做功以外的功}（有用功，如电功）满足
$W_{\text{u}} \ge \Delta G$：$-\Delta G$ 为最大有用功，
即在 $T$ 与 $P$ 固定时可回收的功（燃料电池、
电池所做的功）。
\end{theorem}

\begin{proof}
热力学第一定律：$\Delta U = W + Q$。单一热源下的热力学第二定律：$\Delta S
= Q/T_0 + S_{\text{c}}$，$S_{\text{c}} \ge 0$，故 $Q \le T_0\Delta S$。由此
$W = \Delta U - Q \ge \Delta U - T_0\Delta S = \Delta(U - T_0S) = \Delta F$，利用
两端均有 $T = T_0$。将大气所做功 $-P_0\Delta V$
单独分离出来，$W = W_{\text{u}} - P_0\Delta V$，且 $W_{\text{u}} \ge \Delta U + P_0
\Delta V - T_0\Delta S = \Delta G$。
\end{proof}

\begin{theorem}[演化与平衡判据]\label{thm:b2:thermodynamic-potentials:criteria}
在 $T$、$V$ 固定且不接受功的系统，只能
朝\emph{$F$ 减小}的方向演化；其平衡态是
$F$ 最小的状态。在 $T$、$P$ 固定且不接受有用功的系统，
只能朝\emph{$G$ 减小}的方向演化；其平衡态是
$G$ 最小的状态。（更一般地，对与温度为 $T_0$、压强为 $P_0$ 的外界相接触的系统，
则 $U - T_0S + P_0V$ 减小。）
\end{theorem}

\begin{proof}
在上一定理中取 $W = 0$（或 $W_{\text{u}} = 0$），得 $\Delta F
\le 0$（或 $\Delta G \le 0$），等号对应可逆过程，即
平衡态的演化。
\end{proof}

\begin{example}[等温膨胀：实际与理想情形]\label{ex:b2:thermodynamic-potentials:expansion}
一摩尔气体在 \qty{300}{K} 下由 \qty{1}{L} 膨胀至 \qty{10}{L}，
与热源接触时：$\Delta F = -RT\ln 10 = \qty{-5.7}{kJ}$。
可逆膨胀可回收 \qty{5.7}{kJ}；对抗
压强为 \qty{1}{bar} 的大气膨胀时仅回收 $P_0\Delta V = \qty{0.9}{kJ}$；
自由膨胀入真空则无功可回收。\qty{5.7}{kJ} 即为
气体与其热源在此状态变化中所能给出的最大值，
无论机器多么巧妙。
\end{example}

\section{麦克斯韦关系}

\begin{theorem}[麦克斯韦关系]\label{thm:b2:thermodynamic-potentials:maxwell}
由于 $\dd F$ 与 $\dd G$ 均为全微分（施瓦茨
关于混合二阶偏导数的定理），
\[
\Big(\frac{\partial S}{\partial V}\Big)_T = \Big(\frac{\partial P}{\partial T}\Big)_V,
\qquad
\Big(\frac{\partial S}{\partial P}\Big)_T = -\Big(\frac{\partial V}{\partial T}\Big)_P .
\]
它们将难以直接测量的量（即等温下熵随体积或
压强的变化）通过状态方程表示。由此得到
两个推论：
\[
\Big(\frac{\partial U}{\partial V}\Big)_T = T\Big(\frac{\partial P}{\partial T}\Big)_V - P,
\qquad
C_P - C_V = T\Big(\frac{\partial P}{\partial T}\Big)_V\Big(\frac{\partial V}{\partial T}\Big)_P
= \frac{TV\alpha^2}{\kappa_T} ,
\]
其中 $\alpha = V^{-1}(\partial V/\partial T)_P$ 为膨胀系数，且
$\kappa_T = -V^{-1}(\partial V/\partial P)_T$ 为等温压缩系数。
\end{theorem}

\begin{proof}
按两种次序计算 $\partial^2F/\partial T\partial V$：$-\partial_VS =
-\partial_TP$。对 $G$ 同理。再由 $\dd U = T\dd S - P\dd V$，在 $T$ 固定时
得 $(\partial_VU)_T = T(\partial_VS)_T - P$；而 $C_P - C_V$ 则由
将 $S(T,V(T,P))$ 展开：$(\partial_TS)_P = (\partial_TS)_V +
(\partial_VS)_T(\partial_TV)_P$，再乘以 $T$。
\end{proof}

\begin{example}[关系式的启示]\label{ex:b2:thermodynamic-potentials:reveal}
对理想气体，$T(\partial_TP)_V - P = 0$：内能
与体积无关（焦耳定律得到印证），且 $C_P - C_V = nR$。对范德瓦耳斯
气体，$(\partial_VU)_T = a/V_{\text{m}}^2$：
为吸引作用产生的内压。对液态水，$(\partial_TP)_V = \alpha/\kappa_T = 2.1
\times 10^{-4}/4.6 \times 10^{-10} = \qty{4.6e5}{Pa/K}$：在刚性且装满的容器中加热水
使温度升高一开尔文，压强即升高 \qty{4.6}{bar}
（切勿密封已装满的瓶子），且 $(\partial_VU)_T \approx \qty{1400}{bar}$——
此即液体的内聚力；然而 $c_P - c_V = Tv\alpha^2/\kappa_T \approx
\qty{30}{J/kg/K}$，不足 $1\%$ 的 $c_P$：对凝聚态物质而言，两种
热容几乎相同。
\end{example}

\begin{proposition}[橡胶的熵弹性]\label{prop:b2:thermodynamic-potentials:rubber}
对长度为 $L$ 的橡皮带、受张力 $f$ 作用，$\dd U = T\dd S + f\dd L$
且 $\dd F = -S\dd T + f\dd L$，于是 $(\partial S/\partial L)_T =
-(\partial f/\partial T)_L$，且 $(\partial U/\partial L)_T = f - T(\partial f/
\partial T)_L$。橡胶在较好近似下服从 $f = T\,\varphi(L)$
其中 $\varphi > 0$（当 $L > L_0$ 时）：于是 $(\partial U/\partial L)_T = 0$——即
拉伸时能量不变——且 $(\partial S/\partial L)_T =
-\varphi(L) < 0$：张力是\emph{熵致的}，$f = -T(\partial S/\partial
L)_T$，橡皮带会回拉，因为拉伸后的链段
构型数更少。推论：等温拉伸时它放出热量
（$Q = T\Delta S < 0$）；绝热拉伸时它升温；加热时
在恒定载荷下它收缩（$f/T$ 固定意味着 $\varphi(L)$ 固定，故在
固定 $f$ 时，$\varphi$ 必须随 $T$ 升高而下降）。
\end{proposition}

\begin{proof}
由 $\dd F$ 得麦克斯韦关系，其余为代入。其微观起源是
\cref{ch:b2:particle-diffusion}中的随机游走：
由 $N$ 个长度为 $a$ 的链段组成的链，其端到端距离
服从方差为 $Na^2$ 的高斯分布，故伸长 $L$ 时的构型数
为 $\Omega \propto \eu^{-L^2/2Na^2}$，$S = S_0 -
k_BL^2/2Na^2$，$f = k_BTL/Na^2$：一根劲度系数
\emph{与温度成正比}的弹簧。
\end{proof}

\begin{omfigure}
\begin{tikzpicture}
  \begin{axis}[omaxis, width=6.2cm, height=4.4cm, xmin=0, xmax=400, ymin=0, ymax=1.25,
    xlabel={$T$ (\unit{K})}, ylabel={$f$ （任意单位）}, xtick={100,200,300}, ytick=\empty, clip=false]
    \addplot[omDef, thick, domain=0:380] {x/300};
    \addplot[omDef, thick, domain=0:380] {0.6*x/300};
    \node[font=\scriptsize, omDef, anchor=west] at (axis cs:300,1.05) {$L - L_0$ 较大};
    \node[font=\scriptsize, omDef, anchor=west] at (axis cs:300,0.5) {较小};
    \node[font=\scriptsize] at (axis cs:120,1.1) {橡胶：$f \propto T$};
  \end{axis}
  \begin{axis}[omaxis, width=6.2cm, height=4.4cm, xmin=0, xmax=400, ymin=0, ymax=1.25, xshift=6.6cm,
    xlabel={$T$ (\unit{K})}, ylabel={$f$ （任意单位）}, xtick={100,200,300}, ytick=\empty, clip=false]
    \addplot[omExo, thick, domain=0:380] {1.0-0.0004*x};
    \addplot[omExo, thick, domain=0:380] {0.6-0.00024*x};
    \node[font=\scriptsize] at (axis cs:200,1.15) {钢：$f$ 几乎与 $T$ 无关};
  \end{axis}
\end{tikzpicture}

{\small 固定伸长量下张力随温度的变化。左：橡胶，
$f \propto T$——张力是熵致的，$(\partial f/\partial T)_L > 0$，
拉伸使熵降低，橡皮带升温。右：钢制
弹簧，其张力来自键能，
随金属变软而略有下降。}
\end{omfigure}

\section{两相平衡}

\begin{theorem}[共存条件]\label{thm:b2:thermodynamic-potentials:coexistence}
纯物质在固定 $T$ 与 $P$ 下分为两相，质量分别为
$m_1$、$m_2$，单位质量自由焓分别为 $g_1(T,P)$，
$g_2(T,P)$，则 $G = m_1g_1 + m_2g_2$；平衡时 $G$ 取最小值，于是
\begin{itemize}
\item 若 $g_1 < g_2$，则全部物质进入相 1（稳定相
      是 $g$ \emph{较低}者）；
\item 两相共存仅当
      \[
      g_1(T,P) = g_2(T,P) ,
      \]
      这是 $T$ 与 $P$ 之间的关系，即相图中的
      共存曲线。
\end{itemize}
沿该曲线（克拉珀龙）有
\[
\frac{\dd P}{\dd T} = \frac{s_2 - s_1}{v_2 - v_1} = \frac{L_{1\to2}}{T\,(v_2 - v_1)} .
\]
\end{theorem}

\begin{proof}
$\dd G = (g_1 - g_2)\dd m_1$ 在固定 $T$、$P$（$m_1 + m_2$ 固定）下：
$G$ 随质量向 $g$ 较低的相转移而减小，直至该相
独占全部物质，或直至 $g_1 = g_2$。沿曲线对 $g_1 = g_2$ 求微分，
并对每一相使用 $\dd g = -s\,\dd T + v\,\dd P$：$-s_1\dd T +
v_1\dd P = -s_2\dd T + v_2\dd P$；且 $s_2 - s_1 = L/T$，因为
固定 $T$、$P$ 下的相变是可逆的。本科一年级分册中
用卡诺循环导出了此式；这里它由 $G$ 直接得出。
\end{proof}

\begin{omfigure}
\begin{tikzpicture}
  \begin{axis}[omaxis, width=8.6cm, height=5cm, xmin=0, xmax=10, ymin=-4.6, ymax=1.2,
    xlabel={$T$}, ylabel={$g$}, xtick=\empty, ytick=\empty, clip=false]
    \addplot[omExo, thick, domain=0:10] {0.6-0.18*x};
    \addplot[omDef, thick, domain=0:10] {1.0-0.30*x};
    \addplot[omThm, thick, domain=0:10] {3.0-0.70*x};
    \draw[dashed, gray] (axis cs:3.33,0.9) -- (axis cs:3.33,-1.0);
    \draw[dashed, gray] (axis cs:5.0,0.9) -- (axis cs:5.0,-1.2);
    \node[font=\scriptsize, below] at (axis cs:3.33,-1.0) {$T_{\text{熔}}$};
    \node[font=\scriptsize, below] at (axis cs:5.0,-1.2) {$T_{\text{汽}}$};
    \node[font=\scriptsize, omExo, anchor=west] at (axis cs:8.2,-0.6) {固体};
    \node[font=\scriptsize, omDef, anchor=west] at (axis cs:8.2,-1.6) {液体};
    \node[font=\scriptsize, omThm, anchor=west] at (axis cs:8.2,-3.5) {气体};
    \node[font=\scriptsize] at (axis cs:3.0,-3.9) {斜率 $= -s$：气体最陡};
  \end{axis}
\end{tikzpicture}

{\small 固定压强下三相单位质量的自由焓：各相随温度下降，斜率为 $-s$，
气体最快；稳定相是最低的那条线，相变发生在
两线相交处。相交点之后的上支为
亚稳态
（过冷液体、过热液体）。}
\end{omfigure}

\begin{example}[水：斜率与亚稳态]\label{ex:b2:thermodynamic-potentials:water}
熔化：$L = \qty{334}{kJ/kg}$，$v_{\text{液}} - v_{\text{冰}} = -\qty{9.1e-5}{m^3/kg}$：
$\dd P/\dd T = \qty{-1.3e7}{Pa/K}$，\emph{负}——冰在压强下熔化，
但每开尔文只需 \qty{130}{bar}：滑冰者的 \qty{70}{bar} 压强
只能使熔点降低半度（滑冰靠摩擦力与
预熔表层实现），而位于三
千米冰盖之下的 \qty{270}{bar} 压强使其底部熔化，冰盖得以滑动。汽化在
\qty{100}{\celsius}：$L = \qty{2.26}{MJ/kg}$，$\Delta v = \qty{1.67}{m^3/kg}$：
\qty{3.6}{kPa/K}——山上 \qty{0.7}{bar} 时沸点为 \qty{92}{\celsius}，
高压锅中 \qty{2}{bar} 时为 \qty{120}{\celsius}。低于
\qty{0}{\celsius} 时，液态水满足 $g_{\text{液}} - g_{\text{冰}} \approx
(L/T_{\text{熔}})(T_{\text{熔}} - T) = \qty{1.2}{kJ/kg}$，即每过冷一开尔文
过冷的液态水处于亚稳态而非不稳定态，云中的纯水滴
因缺少凝结核而保持液态，可低至 \qty{-40}{\celsius}。
\end{example}

\begin{omfigure}
\begin{tikzpicture}
  \begin{axis}[omaxis, width=8.4cm, height=5cm, xmin=0.3, xmax=3.2, ymin=0, ymax=1.6,
    xlabel={$v$}, ylabel={$P$}, xtick=\empty, ytick=\empty]
    % vdW-like isotherm below Tc (schematic): P = 8T/(3v-1) - 3/v^2 with T=0.9 in reduced units, v from 0.45..3.2
    \addplot[omDef, thick, domain=0.5:3.2, samples=200] {8*0.9/(3*x-1) - 3/(x*x)};
    \addplot[omExo, thick, dashed, domain=0.45:3.2] {0.647};
    \addplot[draw=none, fill=omExo!25, opacity=0.6] coordinates {(0.603,0.647) (0.603,0.647) (0.616,0.586) (0.628,0.539) (0.640,0.502) (0.652,0.474) (0.664,0.453) (0.676,0.439) (0.689,0.429) (0.701,0.423) (0.713,0.420) (0.725,0.420) (0.737,0.422) (0.750,0.426) (0.762,0.432) (0.774,0.439) (0.786,0.446) (0.798,0.454) (0.810,0.463) (0.823,0.472) (0.835,0.481) (0.847,0.491) (0.859,0.500) (0.871,0.510) (0.884,0.519) (0.896,0.528) (0.908,0.537) (0.920,0.547) (0.932,0.555) (0.944,0.564) (0.957,0.572) (0.969,0.580) (0.981,0.588) (0.993,0.596) (1.005,0.603) (1.017,0.610) (1.030,0.617) (1.042,0.624) (1.054,0.630) (1.066,0.636) (1.078,0.642) (1.091,0.647) (1.091,0.647)};
    \addplot[draw=none, fill=omExo!25, opacity=0.6] coordinates {(1.091,0.647) (1.091,0.647) (1.112,0.656) (1.132,0.664) (1.153,0.672) (1.174,0.678) (1.195,0.685) (1.216,0.690) (1.237,0.695) (1.258,0.700) (1.279,0.704) (1.300,0.708) (1.321,0.711) (1.342,0.714) (1.363,0.716) (1.384,0.718) (1.405,0.720) (1.426,0.721) (1.447,0.722) (1.468,0.723) (1.489,0.724) (1.510,0.724) (1.531,0.724) (1.552,0.724) (1.573,0.724) (1.594,0.723) (1.615,0.722) (1.636,0.722) (1.657,0.721) (1.678,0.719) (1.699,0.718) (1.720,0.717) (1.741,0.715) (1.762,0.714) (1.783,0.712) (1.804,0.710) (1.825,0.708) (1.846,0.706) (1.866,0.704) (1.887,0.702) (1.908,0.700) (1.929,0.698) (1.950,0.696) (1.971,0.693) (1.992,0.691) (2.013,0.688) (2.034,0.686) (2.055,0.684) (2.076,0.681) (2.097,0.679) (2.118,0.676) (2.139,0.673) (2.160,0.671) (2.181,0.668) (2.202,0.666) (2.223,0.663) (2.244,0.660) (2.265,0.658) (2.286,0.655) (2.307,0.652) (2.328,0.650) (2.349,0.647) (2.349,0.647)};
    \node[font=\scriptsize, omExo, anchor=west] at (axis cs:2.3,0.76) {$P_{\text{饱和}}(T)$};
    \node[font=\scriptsize] at (axis cs:0.75,1.35) {液体};
    \node[font=\scriptsize] at (axis cs:2.7,0.35) {蒸气};
    \node[font=\scriptsize] at (axis cs:1.5,0.25) {等面积};
  \end{axis}
\end{tikzpicture}

{\small 低于临界温度的范德瓦耳斯等温线，以及
麦克斯韦作图：真实等温线沿水平线
$P_{\text{饱和}}$ 处的水平线，使两块阴影面积相等——
即条件 $g_{\text{液}} = g_{\text{气}}$，由于 $\dd g = v\,\dd P$ 且
$T$ 固定，回线两端附近上升的部分为
亚稳态的过热液体与过饱和蒸气。}
\end{omfigure}

\begin{remark}[范德瓦耳斯方程与麦克斯韦作图法]\label{rem:b2:thermodynamic-potentials:maxwellconstruction}
范德瓦耳斯方程 $(P + a/v^2)(v - b) = RT$ 在低于
其临界温度 $T_{\text{c}} = 8a/27Rb$ 时给出带回线的等温线。
真实流体并不遵循该回线：在 $P_{\text{饱和}}(T)$ 处它从
由液体支跳到蒸气支，$P_{\text{饱和}}$ 由 $g_{\text{液}}
= g_{\text{气}}$ 确定，等价于沿回线的 $\int_{\text{液}}^{\text{气}}v\,\dd P = 0$：
水平线截取的是\emph{等面积}。回线中
满足 $(\partial P/\partial v)_T > 0$ 的部分为力学不稳定；
它们与直线之间的部分为亚稳态——
即在微波炉中“暴沸”的过热液体与
云室中的过饱和蒸气。
\end{remark}

\begin{method}[使用热力学势]\label{meth:b2:thermodynamic-potentials:method}
(1) 确定约束：孤立 $\to$ $S$ 取极大；$T,V$ 固定 $\to$
$F$ 取极小；$T,P$ 固定 $\to$ $G$ 取极小。(2) 最大功为 $-\Delta F$（
定 $T$ 时）；有用功为 $-\Delta G$（定 $T,P$）。(3) 若需 $S$ 在
定 $T$ 下的导数？用麦克斯韦关系与状态方程。
(4) 两相：比较 $g$；共存条件 $g_1 = g_2$；斜率用克拉珀龙方程，
斜率符号由 $\Delta v$ 的符号决定。(5) 化学
势 $\mu = g$（每摩尔）：物质向 $\mu$ 较低的方向流动。
\end{method}

\section{练习}

\begin{exercise}[$\star$]\label{exo:b2:thermodynamic-potentials:1}
对 $n$ 摩尔理想气体，写出 $F(T,V)$ 与 $G(T,P)$（
至多相差一个只含 $T$ 的函数），由二者分别导出状态方程；证明
由 $S = -\partial_TF$ 可得已知的熵。
\end{exercise}

\begin{exercise}[$\star$]\label{exo:b2:thermodynamic-potentials:2}
一摩尔气体，\qty{300}{K}，在恒温热源中由 \qty{1}{L} 膨胀至 \qty{10}{L}：
求 $\Delta F$、$\Delta U$、$\Delta S$；若过程可逆，求回收的功；
若对抗 \qty{1}{bar} 外压；若向真空膨胀；求
每种情形下产生的熵。
\end{exercise}

\begin{exercise}[$\star$]\label{exo:b2:thermodynamic-potentials:3}
水与冰在 \qty{1}{bar} 下：证明 $g_{\text{液}} - g_{\text{冰}} \approx
(L/T_{\text{熔}})(T_{\text{熔}} - T)$ 在 \qty{0}{\celsius} 附近，并计算
其在 \qty{-5}{} 与 $+\qty{5}{\celsius}$ 处的值（$L = \qty{334}{kJ/kg}$）。哪一相
在各情形下稳定？另一相是否不可能存在？
\end{exercise}

\begin{exercise}[$\star$]\label{exo:b2:thermodynamic-potentials:4}
水的克拉珀龙斜率：熔化（$\Delta v = -\qty{9.1e-5}{m^3/kg}$）与
在 \qty{100}{\celsius} 下的汽化（$L = \qty{2.26}{MJ/kg}$，$\Delta v =
\qty{1.67}{m^3/kg}$）。求 \qty{70}{bar} 下的熔点，以及
在 \qty{0.7}{bar} 与 \qty{2}{bar} 下的沸点。
\end{exercise}

\begin{exercise}[$\star\star$]\label{exo:b2:thermodynamic-potentials:5}
\emph{麦克斯韦关系。}(a) 推导本章的两个麦克斯韦关系
以及由 $\dd U$ 和 $\dd H$ 导出的另两个。(b) 证明 $(\partial U/\partial V)_T =
T(\partial P/\partial T)_V - P$，并对理想气体与
范德瓦耳斯气体求其值。(c) 液态水：$\alpha = \qty{2.1e-4}{K^{-1}}$，$\kappa_T =
\qty{4.6e-10}{Pa^{-1}}$：刚性且装满的容器中，每开尔文的压强升高量，
以及 $(\partial U/\partial V)_T$。(d) 一只装满并密封的玻璃瓶
升温 \qty{10}{K}，请加以评论。
\end{exercise}

\begin{exercise}[$\star\star$]\label{exo:b2:thermodynamic-potentials:6}
\emph{$C_P - C_V$。}(a) 推导 $C_P - C_V = TV\alpha^2/\kappa_T$。(b)
理想气体。(c) 水在 \qty{300}{K} 下；铜（$\alpha = \qty{5e-5}{K^{-1}}$，
$\kappa_T = \qty{7e-12}{Pa^{-1}}$，$\rho = \qty{8900}{kg/m^3}$）；与
$c_P$（\qty{4180}{} 与 \qty{385}{J/kg/K}）比较。(d) 为何可以说固体
有“唯一的”热容？
\end{exercise}

\begin{exercise}[$\star\star$]\label{exo:b2:thermodynamic-potentials:7}
\emph{橡皮筋。}一根橡皮筋服从 $f = CT(L - L_0)$，其中 $C = \qty{1}{N/m/K}$。
(a) 求 $(\partial S/\partial L)_T$ 与 $(\partial U/\partial L)_T$。(b) 求熵变、
热量与功：若在 \qty{300}{K} 下由
$L_0$ 等温拉伸至 $L_0 + \qty{0.1}{m}$。(c) 若改为绝热拉伸（热容
为 \qty{4}{J/K}），求温度变化。(d) 挂上固定重物
并升温 \qty{30}{K}，求 $L - L_0$ 的相对变化。
\end{exercise}

\begin{exercise}[$\star\star$]\label{exo:b2:thermodynamic-potentials:8}
\emph{水洼的蒸发。}水蒸气的化学势
在分压 $p$ 下为 $\mu_{\text{v}} = \mu_{\text{v}}^\circ(T) +
RT\ln(p/p^\circ)$，且在 $p = p_{\text{饱和}}(T)$ 时与液体的相等。
(a) 证明当 $p < p_{\text{饱和}}$ 时液体蒸发，
当 $p > p_{\text{饱和}}$ 时蒸气凝结；定义相对湿度。(b)
求 $\mu_{\text{液}} - \mu_{\text{v}}$，在相对湿度 $50\%$、\qty{293}{K} 下。(c) 为何
洗涤后的衣物即使在 \qty{20}{\celsius} 下也能晾干，而在 $100\%$ 湿度下永远不会？
(d) 露水：它首先在哪些表面形成？为什么？
\end{exercise}

\begin{exercise}[$\star\star$]\label{exo:b2:thermodynamic-potentials:9}
\emph{范德瓦耳斯方程。}$(P + a/v^2)(v - b) = RT$（摩尔量）。(a) 求
临界点（$\partial_vP = \partial_v^2P = 0$）：$v_{\text{c}} = 3b$，
$T_{\text{c}} = 8a/27Rb$，$P_{\text{c}} = a/27b^2$。(b) 对二氧化碳，
$T_{\text{c}} = \qty{304}{K}$，$P_{\text{c}} = \qty{73.8}{bar}$，求 $a$ 与 $b$。(c)
由 $g$ 论证麦克斯韦等面积法则。(d) 回线的哪些部分
是不稳定的，哪些是亚稳态的，以及哪些现象表明
亚稳态的存在？
\end{exercise}

\begin{exercise}[$\star\star\star$]\label{exo:b2:thermodynamic-potentials:10}
\emph{克劳修斯–克拉珀龙方程的积分形式。}对汽化过程，略去 $v_{\text{液}}$
并将蒸气视为理想气体。(a) 证明 $\dd\ln P/\dd T = L_{\text{m}}/RT^2$
（$L_{\text{m}}$ 为摩尔量），并在 $L_{\text{m}}$ 为常量时积分。(b) 水，
$L_{\text{m}} = \qty{40.7}{kJ/mol}$：求 \qty{20}{\celsius} 时的 $P_{\text{饱和}}$，
由 \qty{100}{\celsius} 时的 \qty{1}{bar} 推算，并与实测值 \qty{23}{mbar} 比较。
(c) 预测三相点压强（\qty{0.01}{\celsius}），并与
\qty{611}{Pa} 比较。(d) 证明同一规律可由吉布斯–亥姆霍兹
方程作用于 $\Delta g$ 导出。
\end{exercise}

\begin{exercise}[$\star\star\star$]\label{exo:b2:thermodynamic-potentials:11}
\emph{渗透。}溶剂的化学势因稀溶液中的溶质而降低，
即 $\mu = \mu^\circ - RTx_{\text{s}}$（$x_{\text{s}}$ 为溶质摩尔
分数）。膜只允许溶剂通过；一侧为纯溶剂，
另一侧为溶液，温度同为 $T$。(a) 证明平衡时，
溶液一侧须额外施加压强 $\Pi$，且 $\Pi v_{\text{m}} =
RTx_{\text{s}}$，即 $\Pi = cRT$（范特霍夫）。(b) 海水的 $c \approx
\qty{1.1}{mol/L}$（溶解粒子），温度 \qty{293}{K}：求 $\Pi$。(c) 最小
功，提取一立方米淡水所需；并与实际植物比较
（\qty{3}{kWh/m^3}）。(d) 为什么植物细胞在纯水中会胀破，
而在盐水中会皱缩？
\end{exercise}

\begin{exercise}[$\star\star\star$]\label{exo:b2:thermodynamic-potentials:12}
\emph{链。}一条由 $N$ 个长度为 $a$ 的链节组成的高分子链：
沿带方向、端到端矢量长度为 $L$ 的构型数
$\Omega \propto \eu^{-L^2/2Na^2}$（随机游走）。(a) 求 $S(L)$
并在 $U$ 固定时求张力 $f = -T\,\partial S/\partial L$。(b) 求
该链的刚度及其随温度的变化。(c) 取 $N = 1000$，$a =
\qty{0.5}{nm}$，\qty{300}{K}，$L = \qty{10}{nm}$：求 $f$；需多少根链
并联才能给出 \qty{30}{N}？(d) 估算一种橡胶的杨氏模量，
该橡胶每立方米含 $10^{26}$ 根链，每根链 $N = 100$。
\end{exercise}

\section{问题：橡皮筋与水的相图}

\begin{problem}[{周末问题——会拉伸的熵与相互竞争的
相}]\label{pb:b2:thermodynamic-potentials:1}
\textbf{第一部分——橡皮筋。}
一根质量为 \qty{2}{g}（热容为 \qty{4}{J/K}）的橡皮筋，其张力为
$f = CT(L - L_0)$，其中 $C = \qty{1}{N/m/K}$；在 \qty{300}{K} 且 $L - L_0 =
\qty{0.1}{m}$ 时，$f = \qty{30}{N}$。
\begin{enumerate}
  \item 写出该橡皮筋的 $\dd U$ 与 $\dd F$（$\dd U = T\dd S + f\dd L$）。
  \item 导出麦克斯韦关系 $(\partial S/\partial L)_T = -(\partial
        f/\partial T)_L$.
  \item 证明该橡皮筋满足 $(\partial U/\partial L)_T = 0$，
        这意味着什么？
  \item 在 \qty{300}{K} 下等温拉伸，从 $L_0$ 拉至 $L_0 + \qty{0.1}{m}$：
        求 $\Delta S$、热量 $Q$、功 $W$；并验证 $\Delta U = 0$。
  \item 若改为绝热拉伸，求温度变化。
  \item 橡皮筋悬挂固定重量，将其加热 \qty{30}{K}：求新的
        $L - L_0$。描述“橡皮筋热机”（一个以橡胶为辐条的
        轮子，其一侧的辐条受热）。
  \item 一根由 $N$ 个长度为 $a$ 的链节组成的链，其 $\Omega \propto \eu^{-L^2/2Na^2}$：
        求 $S(L)$ 与 $f(L,T)$；为什么刚度随 $T$ 增大？
  \item 取 $N = 1000$，$a = \qty{0.5}{nm}$，$L = \qty{10}{nm}$，\qty{300}{K}：求
        $f$（单根链），以及产生 \qty{30}{N} 所需的链数。
  \item 为什么钢弹簧的表现恰好相反（固定长度时张力
        随 $T$ 下降）？
\end{enumerate}

\textbf{第二部分——水的相图。}
数据：$L_{\text{熔}} = \qty{334}{kJ/kg}$，$v_{\text{冰}} = \qty{1.091e-3}{m^3/kg}$，
$v_{\text{液}} = \qty{1.000e-3}{m^3/kg}$；在 \qty{100}{\celsius} 时 $L_{\text{气}} = \qty{2.26}{MJ/kg}$，
$v_{\text{气}} = \qty{1.67}{m^3/kg}$；$L_{\text{m}} = \qty{40.7}{kJ/mol}$；
三相点 \qty{0.01}{\celsius}、\qty{611}{Pa}；临界点 \qty{647}{K}、
\qty{221}{bar}。
\begin{enumerate}[resume]
  \item 画出冰、液体与蒸气在 \qty{1}{bar} 下的 $g(T)$ 示意图，并说明
        斜率的次序，并标出各相变点的位置。
  \item 由 $g_1 = g_2$ 导出克拉珀龙公式。
  \item 求熔化曲线的斜率；说明其符号，并解释为何这一情形不寻常。
  \item 一名滑冰者，\qty{700}{N} 的力作用于 \qty{1}{cm^2}：求熔点的变化；
        由此对关于滑冰的流行解释作出结论。
  \item 在 \qty{3}{km} 厚的冰层之下（$\rho = \qty{917}{kg/m^3}$）：求压强与
        熔点；并说明对冰盖的影响。
  \item 求 \qty{100}{\celsius} 时汽化曲线的斜率；求
        \qty{0.7}{bar} 与 \qty{2}{bar} 下的沸点。
  \item 对理想蒸气及常量 $L_{\text{m}}$ 积分克劳修斯–克拉珀龙方程；
        求 \qty{20}{\celsius} 与三相点处的 $P_{\text{饱和}}$；
        并评论其精度。
  \item 在三相点处，比较升华曲线与
        汽化曲线的斜率（$L_{\text{升华}} = L_{\text{熔}} + L_{\text{气}}$）。
  \item 过冷水在 \qty{-10}{\celsius} 时：求 $g_{\text{液}} - g_{\text{冰}}$；
        为什么云滴能保持液态直至 \qty{-40}{\celsius}？
  \item 临界点处，$\Delta v$、$L$ 及斜率会发生什么变化；
        什么是超临界流体？
\end{enumerate}

\textbf{第三部分——应用。}
\begin{enumerate}[resume]
  \item 解释范德瓦耳斯等温线上的麦克斯韦等面积构造，
        依据是 $g$ 的连续性。
  \item 微波炉中过热的水会“爆沸”：位于等温线的哪一支，
        为何危险？
  \item 冷冻干燥：在 $(P,T)$ 图中描述其路径，
        并说明水为何不经熔化便离开。
  \item 为什么一盘水在 \qty{20}{\celsius} 的干燥空气中会蒸发，
        而在饱和空气中却停止——请用化学势说明。
  \item 高压锅内压强为 \qty{2}{bar}：为什么食物烹饪得
        快得多（化学反应速率大约每升高
        \qty{10}{K} 便翻一番）？
  \item 总结：各种约束分别对应哪一个势，
        以及麦克斯韦关系带来了什么。
\end{enumerate}
\end{problem}
